% Vietnamese translation for OI Public Library
% Source-PDF: 图论 GRAPH THEORY/经典算法之二分图 - 何振豪.pdf
\documentclass[11pt]{article}
\usepackage{oipl-vn}

\title{Các thuật toán kinh điển trên đồ thị hai phía}
\author{He Zhenhao}
\date{}

\begin{document}
\maketitle

\origtitle{Classic algorithms for bipartite graphs}
\sourcepath{Graph Theory / Classic Algorithms for Bipartite Graphs - He Zhenhao}

\section{Khái niệm đồ thị hai phía}

Một đồ thị vô hướng được gọi là \term{đồ thị hai phía} nếu có thể chia toàn bộ
tập đỉnh của nó thành hai tập rời nhau $U$ và $V$ sao cho mọi cạnh đều nối một
đỉnh của $U$ với một đỉnh của $V$. Nói cách khác, trong cùng một phía không có
cạnh nội bộ. Cách chia này thường được gọi là một \term{phân hoạch hai phía}
của đồ thị.

Nhiều bài toán ghép cặp trong thi đấu lập trình có dạng tự nhiên là đồ thị hai
phía: một phía là nam, một phía là nữ; một phía là hàng, một phía là cột; hoặc
một phía là bản sao ``đi ra'' của các đỉnh, phía còn lại là bản sao ``đi vào''
của các đỉnh. Vì vậy các thuật toán trên đồ thị hai phía thường là cầu nối từ
một phát biểu tổ hợp sang một thuật toán ghép cặp.

\section{Nhận biết đồ thị hai phía}

Một đồ thị vô hướng $G$ là đồ thị hai phía khi và chỉ khi $G$ không chứa chu
trình độ dài lẻ. Điều kiện này có thể hiểu như sau. Nếu ta đi quanh một chu
trình trong đồ thị hai phía, các đỉnh phải luân phiên giữa $U$ và $V$, nên chỉ
có thể quay về đỉnh xuất phát sau một số chẵn cạnh. Ngược lại, nếu không có
chu trình lẻ, ta có thể tô màu từng thành phần liên thông bằng hai màu mà không
gặp mâu thuẫn.

Phương pháp thường dùng là \term{tô màu hai màu}. Với mỗi thành phần liên thông
chưa xét, chọn một đỉnh bất kỳ, tô nó màu $0$, rồi lan truyền bằng DFS hoặc
BFS:
\begin{enumerate}
  \item Với mỗi cạnh $(u,v)$, nếu $v$ chưa được tô màu thì tô $v$ bằng màu đối
  với màu của $u$.
  \item Nếu $v$ đã được tô màu và màu của $v$ trùng với màu của $u$, đồ thị
  không phải là đồ thị hai phía.
  \item Nếu quá trình kết thúc trên mọi thành phần mà không có mâu thuẫn, hai
  lớp màu chính là hai phía của đồ thị.
\end{enumerate}

Ví dụ điển hình là bài toán chia $n$ học sinh thành hai nhóm sao cho trong cùng
một nhóm không có cặp học sinh quen biết nhau. Mỗi học sinh là một đỉnh, mỗi
quan hệ quen biết là một cạnh. Bài toán hỏi đúng xem đồ thị quan hệ quen biết
có phải là đồ thị hai phía hay không. Lưu ý rằng quan hệ ``$A$ quen $B$'' và
``$B$ quen $C$'' không kéo theo ``$A$ quen $C$''; vì thế ta chỉ dùng các cạnh
được cho trực tiếp.

\section{Ghép cực đại trong đồ thị hai phía}

Một \term{matching} là một tập cạnh đôi một không chung đỉnh. Mục tiêu của
\term{ghép cực đại} là chọn được nhiều cạnh nhất. Trong bài toán tàu lượn, mỗi
toa có hai chỗ và chỉ cho phép một nam ngồi cùng một nữ; mỗi cô gái chỉ chấp
nhận một số chàng trai nhất định. Dựng đồ thị hai phía với một phía là các cô
gái, một phía là các chàng trai, nối cạnh cho mỗi cặp được phép. Số cặp ngồi
được nhiều nhất chính là kích thước ghép cực đại.

\subsection{Đường luân phiên và đường tăng}

Xét một matching hiện tại $M$. Một cạnh thuộc $M$ gọi là cạnh đã ghép; cạnh
không thuộc $M$ gọi là cạnh chưa ghép.

\term{Đường luân phiên} là một đường đi bắt đầu từ một đỉnh chưa ghép và có các
cạnh lần lượt là chưa ghép, đã ghép, chưa ghép, đã ghép, v.v. Một
\term{đường tăng} là đường luân phiên bắt đầu ở một đỉnh chưa ghép và kết thúc
ở một đỉnh chưa ghép, trong đó cạnh đầu và cạnh cuối đều là cạnh chưa ghép.
Trong đồ thị hai phía, nếu ta bắt đầu ở phía trái thì đường tăng sẽ kết thúc ở
phía phải.

Các slide gốc minh họa một đường tăng
\[
  2\to 5\to 1\to 7\to 4\to 8.
\]
Trong hình đó, các cạnh màu đỏ là cạnh đã ghép, các cạnh màu đen là cạnh chưa
ghép. Khi đi theo đường trên, trạng thái cạnh luân phiên giữa chưa ghép và đã
ghép; hai đầu đường là hai đỉnh chưa ghép.

Một đường tăng có các tính chất cơ bản sau:
\begin{enumerate}
  \item Số cạnh trên đường là số lẻ.
  \item Nếu điểm đầu ở phía trái thì điểm cuối ở phía phải.
  \item Các đỉnh trên đường luân phiên giữa hai phía, và các cạnh luân phiên
  giữa chưa ghép và đã ghép.
  \item Chỉ hai đầu đường chưa được ghép; mọi đỉnh trung gian đều đang được
  ghép.
  \item Các cạnh ở vị trí lẻ trên đường là cạnh chưa ghép, còn các cạnh ở vị
  trí chẵn là cạnh đã ghép.
\end{enumerate}

Ý nghĩa của đường tăng nằm ở thao tác đảo trạng thái: biến các cạnh chưa ghép
trên đường thành cạnh đã ghép, đồng thời bỏ các cạnh đã ghép trên đường khỏi
matching. Vì đường tăng có nhiều hơn đúng một cạnh chưa ghép so với cạnh đã
ghép, thao tác này làm kích thước matching tăng thêm $1$.

\subsection{Thuật toán Hungary}

Thuật toán Hungary cho ghép cực đại trong đồ thị hai phía lặp lại đúng ý tưởng
trên:
\begin{enumerate}
  \item Khởi tạo matching $M$ rỗng.
  \item Tìm một đường tăng đối với $M$.
  \item Nếu tìm được, đảo trạng thái các cạnh trên đường tăng để thu được
  matching lớn hơn.
  \item Lặp lại cho đến khi không còn đường tăng.
\end{enumerate}

Định lý Berge bảo đảm rằng một matching là cực đại khi và chỉ khi không tồn tại
đường tăng đối với nó. Do đó vòng lặp trên dừng lại với một ghép cực đại.

Trong cài đặt thường gặp, ta duyệt từng đỉnh phía trái và dùng DFS để thử tìm
đường tăng bắt đầu từ đỉnh đó. Mỗi lần DFS có thể quét qua $O(E)$ cạnh, và có
$O(V)$ lần thử, nên độ phức tạp là $O(VE)$. Có thể dùng BFS để tổ chức quá
trình tìm kiếm, nhưng lõi tư tưởng vẫn là tìm đường tăng.

\subsection{Thuật toán Hopcroft--Karp}

Nếu đồ thị có hàng nghìn đỉnh và khá dày, việc mỗi lần chỉ tăng matching thêm
một cạnh có thể chậm. Thuật toán \term{Hopcroft--Karp} cải thiện bằng cách tìm
nhiều đường tăng ngắn nhất, đôi một không giao nhau theo đỉnh, trong cùng một
pha.

Một pha của Hopcroft--Karp gồm hai bước:
\begin{enumerate}
  \item BFS từ tất cả đỉnh chưa ghép ở phía trái để xây các lớp khoảng cách
  trong đồ thị luân phiên. Việc này xác định độ dài nhỏ nhất của một đường
  tăng hiện tại.
  \item DFS theo các lớp vừa xây để lấy một họ cực đại các đường tăng ngắn
  nhất không giao nhau theo đỉnh, rồi đảo trạng thái đồng thời trên các đường
  này.
\end{enumerate}

Có thể chứng minh số pha chỉ là $O(\sqrt{|V|})$, vì thế độ phức tạp tổng quát
là
\[
  O(|E|\sqrt{|V|}).
\]
Các phần sau đều dựa trên cùng hạt nhân này: đưa bài toán về ghép trong đồ thị
hai phía, rồi dùng đường tăng hoặc một thuật toán ghép nhanh hơn để giải.

\section{Phủ đỉnh nhỏ nhất trong đồ thị hai phía}

Một \term{phủ đỉnh} là một tập đỉnh sao cho mỗi cạnh của đồ thị có ít nhất một
đầu mút nằm trong tập đó. Bài toán \term{phủ đỉnh nhỏ nhất} yêu cầu chọn ít
đỉnh nhất có thể.

Định lý König phát biểu rằng trong đồ thị hai phía,
\[
  \text{kích thước ghép cực đại}
  =
  \text{kích thước phủ đỉnh nhỏ nhất}.
\]
Vì vậy, nếu chỉ cần kích thước phủ nhỏ nhất, ta có thể tính ghép cực đại.

Một ví dụ trong slide là bài toán bắn laser trên lưới $N\times M$. Một số ô có
vật thể cần loại bỏ. Laser được đặt ở biên, mỗi lần bắn chỉ đi theo một hàng
hoặc một cột. Hỏi cần ít nhất bao nhiêu lần bắn để mọi vật thể đều bị bắn
trúng.

Dựng đồ thị hai phía như sau: phía trái gồm các hàng $1,\ldots,N$, phía phải
gồm các cột $1,\ldots,M$. Với mỗi vật thể ở ô $(i,j)$, thêm cạnh giữa hàng $i$
và cột $j$. Chọn bắn một hàng hoặc một cột chính là chọn một đỉnh. Mọi vật thể
bị bắn trúng khi và chỉ khi mọi cạnh được phủ bởi ít nhất một đỉnh đã chọn.
Do đó đáp án là kích thước phủ đỉnh nhỏ nhất, và theo định lý König bằng kích
thước ghép cực đại của đồ thị vừa dựng. Trong hình minh họa của slide, có bốn
vật thể nhưng chỉ cần chọn hai đường thẳng là đủ.

Nếu cần khôi phục chính tập đỉnh phủ nhỏ nhất, sau khi có một ghép cực đại ta
có thể làm chuẩn như sau. Từ các đỉnh chưa ghép ở phía trái, đi theo cạnh chưa
ghép từ trái sang phải và cạnh đã ghép từ phải sang trái, đánh dấu tập đỉnh
đạt được là $Z$. Khi đó một phủ đỉnh nhỏ nhất là
\[
  (U\setminus Z)\cup (V\cap Z).
\]

\section{Phủ đường đi nhỏ nhất}

Bài toán \term{phủ đường đi nhỏ nhất} yêu cầu chọn ít nhất số đường đi đơn,
đôi một không giao nhau theo đỉnh, sao cho mỗi đỉnh của đồ thị được nằm trong
đúng một đường đi. Trong ngữ cảnh chuẩn của công thức dưới đây, đồ thị ban đầu
là một DAG.

Với DAG $G=(V,E)$, tách mỗi đỉnh $x$ thành hai bản sao $x_L$ và $x_R$, rồi
dựng đồ thị hai phía
\[
  U=\{x_L\mid x\in V\},\qquad
  W=\{x_R\mid x\in V\}.
\]
Với mỗi cung $x\to y$ của DAG, thêm cạnh $x_Ly_R$ vào đồ thị hai phía. Khi
chọn một cạnh trong matching, ta hiểu là nối $x$ đứng trước $y$ trên một đường
đi. Vì matching không dùng chung đỉnh, mỗi đỉnh trong DAG có nhiều nhất một
cung được chọn đi ra và nhiều nhất một cung được chọn đi vào; do đó các cạnh
được chọn tạo thành một họ đường đi rời nhau.

Ban đầu nếu không chọn cạnh nào, ta có $|V|$ đường đi đơn lẻ. Mỗi cạnh ghép
được chọn nối hai đường đi lại với nhau và làm số đường giảm đi $1$. Vì vậy:
\[
  \text{phủ đường đi nhỏ nhất}
  =
  |V|-\text{kích thước ghép cực đại}.
\]

\section{Tập độc lập lớn nhất trong đồ thị hai phía}

Một \term{tập độc lập} là một tập đỉnh sao cho không có cạnh nào nối hai đỉnh
cùng nằm trong tập đó. Bài toán hỏi số đỉnh nhiều nhất có thể chọn.

Phủ đỉnh và tập độc lập là hai khái niệm bù nhau trong mọi đồ thị. Nếu $C$ là
một phủ đỉnh, thì $V\setminus C$ là tập độc lập: nếu tồn tại một cạnh có cả hai
đầu mút trong $V\setminus C$, cạnh đó không được $C$ phủ. Ngược lại, nếu $I$
là tập độc lập, thì $V\setminus I$ là phủ đỉnh: mỗi cạnh phải có ít nhất một
đầu mút không thuộc $I$.

Do đó trong đồ thị hai phía,
\[
  \alpha(G)
  =
  |V|-\tau(G)
  =
  |V|-\nu(G),
\]
trong đó $\alpha(G)$ là kích thước tập độc lập lớn nhất, $\tau(G)$ là kích
thước phủ đỉnh nhỏ nhất, và $\nu(G)$ là kích thước ghép cực đại. Nói cách
khác, với đồ thị hai phía, đáp án tập độc lập lớn nhất bằng tổng số đỉnh trừ
đi kích thước ghép cực đại.

\section{Ghép hoàn hảo tối ưu trong đồ thị hai phía}

Xét một đồ thị hai phía đầy đủ với hai phía $X$ và $Y$ có cùng kích thước $n$.
Mỗi cạnh $(i,j)$ có trọng số $w_{ij}$, trọng số có thể âm. Bài toán
\term{ghép hoàn hảo tối ưu} yêu cầu chọn một ghép hoàn hảo có tổng trọng số
lớn nhất.

Thuật toán kinh điển cho bài toán cực đại này là \term{Kuhn--Munkres}, thường
gọi tắt là KM. Hai khái niệm chính là nhãn khả thi và đồ thị bằng nhau.

Một hệ \term{nhãn khả thi} gồm các giá trị $L_x(i)$ với $i\in X$ và $L_y(j)$
với $j\in Y$ sao cho
\[
  L_x(i)+L_y(j)\ge w_{ij}
  \qquad\text{với mọi }i,j.
\]
\term{Đồ thị bằng nhau} là đồ thị con gồm đúng các cạnh thỏa
\[
  L_x(i)+L_y(j)=w_{ij}.
\]

Nếu đồ thị bằng nhau có một ghép hoàn hảo, thì ghép đó là tối ưu trong đồ thị
ban đầu. Thật vậy, với mọi ghép hoàn hảo $M$ ta có
\[
  \sum_{(i,j)\in M} w_{ij}
  \le
  \sum_{(i,j)\in M} \bigl(L_x(i)+L_y(j)\bigr)
  =
  \sum_i L_x(i)+\sum_j L_y(j).
\]
Nếu $M$ nằm hoàn toàn trong đồ thị bằng nhau, bất đẳng thức trở thành đẳng
thức, nên không ghép hoàn hảo nào có tổng trọng số lớn hơn $M$.

Khung thuật toán KM như sau:
\begin{enumerate}
  \item Khởi tạo nhãn khả thi, thường đặt
  \[
    L_x(i)=\max_j w_{ij},\qquad L_y(j)=0.
  \]
  \item Trên đồ thị bằng nhau hiện tại, dùng thuật toán Hungary để tìm ghép.
  \item Nếu đã có ghép hoàn hảo thì dừng.
  \item Nếu chưa, điều chỉnh nhãn để vẫn giữ điều kiện
  $L_x(i)+L_y(j)\ge w_{ij}$ nhưng làm đồ thị bằng nhau có thêm cạnh, rồi quay
  lại bước tìm ghép.
\end{enumerate}

Trong cách cài đặt chuẩn $O(n^3)$, khi một lần tìm đường tăng bị kẹt, ta xét
tập đỉnh $S\subseteq X$ và $T\subseteq Y$ đã được thăm trong cây luân phiên,
rồi lấy
\[
  \Delta=\min_{i\in S,\ j\notin T}
  \bigl(L_x(i)+L_y(j)-w_{ij}\bigr).
\]
Sau đó giảm $L_x(i)$ đi $\Delta$ với mọi $i\in S$ và tăng $L_y(j)$ thêm
$\Delta$ với mọi $j\in T$. Các cạnh đang bằng nhau trong cây vẫn bằng nhau,
điều kiện khả thi vẫn đúng, và ít nhất một cạnh mới từ $S$ sang $Y\setminus T$
trở thành cạnh bằng nhau. Lặp lại quá trình này cho đến khi thu được ghép hoàn
hảo.

Nếu bài toán yêu cầu tổng trọng số nhỏ nhất, có thể đổi dấu mọi trọng số rồi
chạy thuật toán cực đại. Nếu đồ thị không đầy đủ, ta có thể thêm các cạnh giả
với trọng số rất nhỏ cho bài toán cực đại, hoặc rất lớn cho bài toán cực tiểu,
miễn là cách thêm cạnh không làm sai ý nghĩa nghiệm hợp lệ.

\section{Ghi chú về đồ thị tổng quát}

Các thuật toán ở trên khai thác mạnh cấu trúc hai phía: cạnh luôn đi từ trái
sang phải, đường tăng có dạng luân phiên rõ ràng, và các định lý như König
được dùng trực tiếp. Nếu bài toán là ghép cực đại trong đồ thị tổng quát, ta
không thể dùng nguyên các thuật toán này; thuật toán kinh điển khi đó là
thuật toán Blossom.

Các slide gốc có nhắc đến mã nguồn đính kèm cho Hungary, Hopcroft--Karp và KM,
nhưng phần đính kèm không nằm trong nội dung văn bản trích xuất từ PDF. Vì vậy
bản ghi chú này chỉ trình bày ý tưởng, định lý và khung cài đặt tự chứa.

\end{document}
